\Needspace{12\baselineskip}
\section{References}

\begingroup
\setlength{\parskip}{0pt}
\titlespacing*{\subsection}{0pt}{1.0ex}{0.3ex}

\subsection*{Sacred Scripture}

{\fontsize{8.6}{9.6}\selectfont
\begin{itemize}[itemsep=0pt,parsep=0pt,topsep=0.05em]
\item Wisdom 13:1--9 and Romans 1:19--20, quoted from the Douay--Rheims translation (Challoner revision; public domain); working web witness checked 2026-07-24 at \sourceurl{https://drbo.org/chapter/25013.htm}{drbo.org (Wisdom 13)} and \sourceurl{https://drbo.org/chapter/52001.htm}{drbo.org (Romans 1)}.
\end{itemize}}

\subsection*{Authoritative Catholic witnesses}

{\fontsize{8.6}{9.6}\selectfont
\begin{itemize}[itemsep=0pt,parsep=0pt,topsep=0.05em]
\item First Vatican Council, dogmatic constitution \emph{Dei Filius} (24 April 1870), ch.~2 \emph{De revelatione} with canon 1 on revelation; canons on God no.~2 (materialism); and ch.~4 \emph{De fide et ratione} (the two orders, the mutual-aid sentence, the just liberty of the sciences); Holy See Latin text, \sourceurl{https://www.vatican.va/content/pius-ix/la/documents/constitutio-dogmatica-dei-filius-24-aprilis-1870.html}{vatican.va}, checked 2026-07-24 and 2026-07-25. English glosses of the canons and ch.~4 passages are labeled working glosses by this essay; the Latin governs.
\item \emph{Catechism of the Catholic Church} 31--35 (ways of coming to know God; ``converging and convincing arguments''), \sourceurl{https://www.vatican.va/content/catechism/en/part_one/section_one/chapter_one/ii_ways_of_coming_to_know_god.html}{official English}; 36--38 (the knowledge of God according to the Church, quoting \emph{Dei Filius} and \emph{Humani generis}), \sourceurl{https://www.vatican.va/content/catechism/en/part_one/section_one/chapter_one/iii_the_knowledge_of_god_according_to_the_church.html}{official English}; 39 (speaking of God with all men), \sourceurl{https://www.vatican.va/content/catechism/en/part_one/section_one/chapter_one/iv_how_can_we_speak_about_god.html}{official English}. All checked 2026-07-24.
\item John Paul~II, encyclical \emph{Fides et ratio} (14 September 1998), nos.~67 (fundamental theology and the naturally knowable truths; ``truly propaedeutic path to faith''), 83 (philosophy of genuinely metaphysical range), and 88 (scientism); \sourceurl{https://www.vatican.va/content/john-paul-ii/en/encyclicals/documents/hf_jp-ii_enc_14091998_fides-et-ratio.html}{official English}, checked 2026-07-24 (nos.~83, 88) and 2026-07-25 (no.~67).
\end{itemize}}

\subsection*{Thomistic framework}

{\fontsize{8.6}{9.6}\selectfont
\begin{itemize}[itemsep=0pt,parsep=0pt,topsep=0.05em]
\item Thomas Aquinas, \emph{Summa theologiae} I, q.~2, a.~2 (demonstration \emph{quia} of God's existence from effects; ad~1, the \emph{praeambula fidei}); I, q.~2, a.~3 (objections 1--2 and the opening of the corpus, quoted for the tradition's method; the five ways themselves outside scope); and I, q.~12, a.~12 (sense-origin, reach, and limits of natural knowledge of God). Latin checked 2026-07-24 and 2026-07-25 in the \sourceurl{https://www.corpusthomisticum.org/iopera.html}{Corpus Thomisticum} online corpus; English from the public-domain translation of the Fathers of the English Dominican Province via \sourceurl{https://www.newadvent.org/summa/1002.htm}{New Advent (q.~2)} and \sourceurl{https://www.newadvent.org/summa/1012.htm}{New Advent (q.~12)}. Work, question, article, and division govern over either site's presentation.
\end{itemize}}

\subsection*{Positivist and empiricist primary texts}

{\fontsize{8.6}{9.6}\selectfont
\begin{itemize}[itemsep=0pt,parsep=0pt,topsep=0.05em]
\item David Hume, \emph{An Enquiry Concerning Human Understanding} (1748), sect.~IV (induction) and sect.~XII, part~III (the closing paragraph, E 12.34; SBN 165); public domain; critical web text checked 2026-07-24 at \sourceurl{https://davidhume.org/texts/e/12}{davidhume.org}.
\item Auguste Comte, \emph{Cours de philosophie positive} (1830--1842), as condensed and translated by Harriet Martineau, \emph{The Positive Philosophy of Auguste Comte}, vol.~1 (London: George Bell \& Sons, 1896): pp.~1--4 (the law of the three stages, its grounds, and the theory--fact circle), pp.~15--17 (intellectual anarchy and social reorganization), pp.~148--149 (the astronomical limits of the knowable); public domain; scan of the 1896 printing checked 2026-07-24 and 2026-07-25 via the \sourceurl{https://archive.org/details/dli.ernet.524059}{Internet Archive}.
\item \emph{Wissenschaftliche Weltauffassung: Der Wiener Kreis} (Vienna: Artur Wolf Verlag, 1929; Ver\"offentlichungen des Vereines Ernst Mach; 59 pp.; no authors named on the title page, foreword signed for the Verein by Hans Hahn, Otto Neurath, and Rudolf Carnap). German text (1929) in the public domain in the United States; quoted in German with labeled working glosses by this essay. Wording checked 2026-07-25 at the transcription of the \sourceurl{https://www.phil.cmu.edu/projects/carnap/editorial/latex_pdf/1929-5.pdf}{Collected Works of Rudolf Carnap editorial project} (phil.cmu.edu); the 1929 printing governs.
\item A.~J. Ayer, \emph{Language, Truth and Logic} (London: Victor Gollancz, 1936; 2nd ed.\ 1946), preface, ch.~1 (the verifiability criterion; strong and weak senses; rejection of conclusive verifiability and falsifiability), ch.~6 (theology), and the 1946 introduction (the retreat: strong/weak collapse, the ``Absolute is lazy'' concession, direct/indirect verifiability, the criterion as definition); in copyright, quoted briefly; wording checked 2026-07-24 and 2026-07-25 in an \sourceurl{https://archive.org/details/AlfredAyer}{Internet Archive scan} of the Pelican reprint, which prints the 1946 introduction as an appendix.
\item W.~V. Quine, ``Two Dogmas of Empiricism,'' \emph{Philosophical Review} 60 (1951): 20--43; revised in \emph{From a Logical Point of View} (Cambridge, Mass.: Harvard University Press, 1953). In copyright; quoted minimally under fair-quotation practice. Wording checked 2026-07-25 at the web transcription edited by Andrew Chrucky, \sourceurl{https://www.ditext.com/quine/quine.html}{ditext.com}, which marks the 1951 and 1961 readings; the transcription is a wording-check witness only, and the printed text governs.
\end{itemize}}

\subsection*{Demarcation and philosophy-of-science scholarship}

{\fontsize{8.6}{9.6}\selectfont
\begin{itemize}[itemsep=0pt,parsep=0pt,topsep=0.05em]
\item Karl R. Popper, \emph{Logik der Forschung} (Vienna: Springer, imprinted 1935, issued 1934), translated as \emph{The Logic of Scientific Discovery} (London: Hutchinson, 1959); cited bibliographically, not collated. ``Science: Conjectures and Refutations'' (1953 lecture), ch.~1 of \emph{Conjectures and Refutations: The Growth of Scientific Knowledge} (London: Routledge \& Kegan Paul, 1963); in copyright, quoted briefly; wording checked 2026-07-25 at an unofficial \sourceurl{https://emilkirkegaard.dk/en/wp-content/uploads/Karl-Popper-SCIENCE-CONJECTURES-AND-REFUTATIONS.pdf}{web transcription} used as a wording-check witness only; no printed edition was collated, and the printed text governs.
\item Carl G. Hempel, ``Problems and Changes in the Empiricist Criterion of Meaning,'' \emph{Revue Internationale de Philosophie} 4, no.~11 (1950): 41--63; reprinted in Leonard Linsky, ed., \emph{Semantics and the Philosophy of Language} (Urbana: University of Illinois Press, 1952), 163--185; revised as ``Empiricist Criteria of Cognitive Significance: Problems and Changes'' in \emph{Aspects of Scientific Explanation} (New York: Free Press, 1965). Identified but not collated; no Hempel wording is quoted anywhere in this essay; the account of the survey's structure and conclusion is used as reported in the Fetzer witness below.
\item Isaiah Berlin, ``Verification,'' \emph{Proceedings of the Aristotelian Society} 39 (1938--39): 225--248 (the objection Ayer credits in 1946 as ``Verifiability in Principle''), and Alonzo Church, review of Ayer's \emph{Language, Truth and Logic}, 2nd ed., \emph{Journal of Symbolic Logic} 14 (1949): 52--53; cited from standard bibliographic records as the classic technical objections; not collated. Berlin's result is additionally witnessed by Ayer's own 1946 concession, quoted from the Ayer scan above.
\item Rudolf Carnap, \emph{The Logical Syntax of Language} (1934; English trans.\ London: Kegan Paul, 1937); cited bibliographically, not collated, for the principle of tolerance and the ascent of philosophy to the logic of science.
\item James Fetzer, ``Carl Hempel,'' \emph{Stanford Encyclopedia of Philosophy} (substantive revision 16 March 2026), \sourceurl{https://plato.stanford.edu/entries/hempel/}{plato.stanford.edu}, checked 2026-07-24 and 2026-07-25, for the standard account of the criterion's successive formulations and failure.
\item Thomas Uebel, ``The Vienna Circle,'' \emph{Stanford Encyclopedia of Philosophy} (substantive revision 7 May 2024), \sourceurl{https://plato.stanford.edu/entries/vienna-circle/}{plato.stanford.edu}, checked 2026-07-25, for the Circle's membership, the manifesto's attribution, and the dispersal of the 1930s.
\item Stephen Thornton, ``Karl Popper,'' \emph{Stanford Encyclopedia of Philosophy} (substantive revision 12 September 2022), \sourceurl{https://plato.stanford.edu/entries/popper/}{plato.stanford.edu}, checked 2026-07-25, for the publication history of \emph{Logik der Forschung} and the standard account and criticisms of the demarcation criterion.
\end{itemize}}

\subsection*{Contemporary scientism exhibits}

{\fontsize{8.6}{9.6}\selectfont
\begin{itemize}[itemsep=0pt,parsep=0pt,topsep=0.05em]
\item Stephen Hawking and Leonard Mlodinow, \emph{The Grand Design} (New York: Bantam, 2010), p.~5 (``philosophy is dead''); in copyright, not collated; the sentence and page locus checked 2026-07-25 at the named study-page witness of C.~R. Nave, \sourceurl{http://hyperphysics.gsu.edu/Nave-html/Faithpathh/GrandDesign.html}{HyperPhysics (Georgia State University)}.
\item Lawrence M. Krauss, ``The Consolation of Philosophy,'' \emph{Scientific American} (online), 27 April 2012; in copyright, quoted briefly; wording checked 2026-07-25 at \sourceurl{https://www.scientificamerican.com/article/the-consolation-of-philos/}{scientificamerican.com}.
\item Alex Rosenberg, \emph{The Atheist's Guide to Reality: Enjoying Life without Illusions} (New York: W.~W. Norton, 2011), p.~6 and the opening-chapter list of ``persistent questions''; in copyright, not collated; quoted as reported 2026-07-25 at the named review witness, Neil Shenvi, ``A Short Review of Rosenberg's \emph{The Atheist's Guide to Reality},'' \sourceurl{https://shenviapologetics.com/a-short-review-of-rosenbergs-the-atheists-guide-to-reality/}{shenviapologetics.com}, corroborated across the book's critical literature.
\end{itemize}}

\endgroup
